% Body-local compatibility definitions; the official Template.tex header is unchanged.
\providecommand{\toprule}{\hline}
\providecommand{\midrule}{\hline}
\providecommand{\bottomrule}{\hline}

\ninept
\maketitle

\begin{abstract}
Where perturbation lands can matter as much as how much is added. We study
whether mask policy alone can redirect release error toward motion-salient
regions when the release backend is fixed. We introduce region-adaptive token
(RAT) fusion, which combines VideoMAE semantic salience, temporal variation,
Sobel texture, and RAFT flow to allocate block-level perturbations, together
with a same-backend mask-isolation protocol. On a fixed balanced subset of 20
UCF101 videos, RAT beats Random and Shifted fusion on Cover, RawErr, and IncErr
for all 20/20 videos at \(k=20\), with RAT-minus-Random differences of 0.1520,
0.1360, and 0.1486 and RAT-minus-Shifted differences of 0.1718, 0.1522, and
0.1649. The advantage persists across ten master seeds (1,400 successful runs)
and \(k\in\{10,20,30,40\}\), while RAT+traj. is statistically
indistinguishable from RAT. These controlled comparisons show that spatial
alignment, rather than score concentration alone, drives the gain, establishing
mask policy as a controllable design lever for motion-salience-aligned
perturbation allocation. For fixed or external masks, the selected low-frequency
DCT queries additionally admit mask-restricted R\'enyi-DP accounting.
\end{abstract}

\begin{keywords}
Video privacy, perturbation allocation, spatiotemporal masking, motion
salience, R\'enyi differential privacy
\end{keywords}

\section{Introduction}
\label{sec:intro}

Where perturbation lands can matter as much as how much perturbation is added.
Motion-salient regions move across frames, so a fixed spatial region can miss
the signal that carries motion and activity. Uniform perturbation spends error
on broad areas, while detector-based anonymization couples localization to
predefined object categories. More generally, existing video-privacy pipelines
often couple masking or localization with anonymization, the release backend,
or recognition evaluation. That coupling makes the independent effect of mask
allocation difficult to isolate.

We isolate the central allocation question: with the backend input \(Z\),
selected area \(k\), DCT-Gaussian mechanism \(\mathcal{G}\), calibration
\(\sigma\), accounting target, and mechanism RNG \(\xi\) controlled, can
changing only \(M^\pi\) systematically redirect distortion toward
motion-salient regions? Equation (2) makes this design explicit: the output
changes through the mask sequence \(M^\pi\), which is the experimental
variable. RAT cue fusion supplies the proposed mask policy, and the comparisons
include Random, Shifted fusion, VideoMAE-only, RAT+traj., and a proxy-aligned
Motion reference.

The study makes three primary contributions. First, it formulates perturbation
placement as an independently testable design variable and introduces a
same-backend mask-isolation protocol. Second, RAT fuses semantic, motion,
texture, and temporal cues for block-level perturbation allocation. Third,
controlled evidence spans videos, stochastic realizations, selected-area
budgets, and ablations, identifying the role of spatial alignment and cue
fusion. Fixed-mask selected-query R\'enyi-DP accounting provides a technical
foundation for the release backend and its formal scope.

\section{Related Work}
\label{sec:related}

\noindent\textbf{(a) Video privacy and anonymization.}
MaskDP studies activity recognition on avatar-anonymized datasets with masked
differential privacy~\cite{schneider2024maskdp}, and Video-DPRP studies a
differentially private approach for privacy-preserving action
recognition~\cite{nken2025videodprp}. STPrivacy treats video as tubelets for
video-level privacy-preserving action recognition, using adaptive token
selection and anonymization~\cite{li2023stprivacy}. Ilic et
al.~\cite{ilic2024selective} and Aslam et al.~\cite{aslam2026tokenprune}
address related selective or temporally consistent anonymization settings. These
works combine masks or transformations with recognition objectives and
evaluation pipelines. They therefore study complete systems, whereas this work
isolates the allocation variable by fixing the DCT-Gaussian backend and
calibration and changing only the mask.

\noindent\textbf{(b) Motion-aware masking.}
Motion-guided masking uses motion cues for spatiotemporal representation
learning~\cite{fan2023mgm}. VideoMAE supplies the semantic-token representation
used as the dominant RAT cue~\cite{tong2022videomae}, and RAFT supplies the
auxiliary flow cue~\cite{teed2020raft}. We use these components as scoring
signals, repurposing these representations as complementary allocation cues
and directly testing their value under a controlled release backend. The
experiment compares semantic, temporal, texture, and flow cue fusion with
random, shifted, and single-cue controls.

\noindent\textbf{(c) Gaussian and R\'enyi-DP accounting.}
Differential privacy calibrates randomization to query
sensitivity~\cite{dwork2006dp}, and R\'enyi-DP composes repeated Gaussian
releases~\cite{mironov2017rdp}. These principles provide the selected-query
accounting foundation for the fixed-mask backend, while the experiments focus
on the independent allocation question.

\section{Method}
\label{sec:method}

\begin{figure*}[t]
\centering
  \includegraphics[width=0.98\textwidth]{figures/rat_method_overview.pdf}
\caption{Overview of the proposed motion-salience-aligned perturbation
allocation pipeline. The raw video provides a common backend input and
multi-cue salience scores; the top-$k$ mask selects blocks for the fixed-mask
DCT-Gaussian release, while unselected blocks are copied from the backend
input.}
\label{fig:main}
\vspace{-0.25in}
\end{figure*}

\subsection{Perturbation-allocation setup}

Let \(X=(x_1,\ldots,x_T)\) denote the raw video used for evaluation and
\(Z=(z_1,\ldots,z_T)\) the common backend input to the controlled
DCT-Gaussian mechanism. Each frame is partitioned into a \(B\times B\) block
grid, with \(B=14\) in our implementation. A mask policy \(\pi\) produces
scores \(S_t^\pi\in\mathrm{R}^{B\times B}\) and selects
\begin{equation}
M_t^\pi=\operatorname{TopK}_k(S_t^\pi), \qquad |M_t^\pi|=k .
\end{equation}
Equation (2) is the central controlled design: for every clip and policy, all
terms except the mask sequence \(M^\pi\) are fixed. The backend input \(Z\),
selected-area budget \(k=20\), DCT-Gaussian mechanism \(\mathcal{G}\), Gaussian
calibration \(\sigma\), accounting target, and Gaussian RNG stream \(\xi\) are
matched across policies:
\begin{equation}
Y^\pi=\mathcal{G}(Z,M^\pi;\sigma,\xi).
\end{equation}
Thus, \(M^\pi\) is the experimental variable. In the saved runs, \(Z\) is the
common pre-release clip-preprocessing and flow-guided merge output from the
original RAT-DP directory; all policies reuse it. The Gaussian calibration and
RNG stream are matched across policies, while selected-block ordering remains
policy-specific.

To evaluate allocation, let \(R_t\) be a top-\(k\) motion-proxy mask computed
from frame differences in the raw video, and write \(R(X)=\{R_t\}_{t=1}^T\).
For any nonnegative per-pixel distortion map \(D\), define
\begin{equation}
\Phi_R(D)=
\frac{\sum_{t,p\in R_t}D_{t,p}}{\sum_{t,p}D_{t,p}} .
\end{equation}
We report mask-proxy overlap, raw-video release error share, and incremental
backend perturbation share:
\begin{equation}
\begin{gathered}
A_{\mathrm{cov}}(\pi)=
\frac{\sum_t |M_t^\pi\cap R_t|}{\sum_t |R_t|},\\
A_{\mathrm{raw}}(\pi)=\Phi_R(D^X), \qquad
D^X_{t,p}=\|X_t(p)-Y_t^\pi(p)\|_2^2,\\
A_{\mathrm{inc}}(\pi)=\Phi_R(D^Z), \qquad
D^Z_{t,p}=\|Z_t(p)-Y_t^\pi(p)\|_2^2 .
\end{gathered}
\end{equation}
RawErr therefore measures total release error relative to the raw evaluation
video, while IncErr isolates distortion newly introduced by the fixed
DCT-Gaussian backend relative to \(Z\). The motion proxy is an operational
reference for quantifying motion-aligned allocation, and these quantities are
empirical allocation scores.

\subsection{RAT cue-fusion policy}

RAT computes a block score from four normalized cues:
\begin{equation}
s_{t,b}=0.85\tilde a_{t,b}+0.10\tilde d_{t,b}
       +0.03\tilde g_{t,b}+0.02\tilde m_{t,b}.
\end{equation}
Here \(\tilde a\) is VideoMAE semantic-token salience~\cite{tong2022videomae},
\(\tilde d\) is tubelet temporal variation, \(\tilde g\) is Sobel texture, and
\(\tilde m\) is RAFT flow magnitude~\cite{teed2020raft}. Each cue is min-max
normalized on the block grid and the fused map is spatially smoothed. These
weights are fixed a priori and are not tuned against the reported evaluation
proxy, avoiding metric-targeted tuning.
All video-dependent cues for RAT fusion, RAT+traj., and VideoMAE-only are
computed from \(X\) after the same frame resizing and clip preprocessing used
by the saved run; \(Z\) is used as the common release-backend input, not to
compute these selector scores.

The main RAT-fusion row selects the top-\(k\) blocks from this per-frame fused
score. We also evaluate a trajectory-smoothed variant, RAT+traj., that
propagates the same scores along block trajectories induced by average RAFT
flow over tubelets of length 2. For a trajectory \(j\), the path score is
\begin{equation}
u_j=0.6\,\mathrm{mean}(s_j)+0.4\,\mathrm{mean}(\mathrm{top25\%}(s_j)).
\end{equation}
We add selector noise with standard deviation \(0.05\), apply a softmax with
inverse temperature 8, and update frame scores with a flow-confidence decay
\(\exp(-\|v\|/12)\) and temporal smoothing weight 0.8. The top-\(k\) blocks
of the resulting score map form the RAT+traj. mask. This selector noise is
distinct from Gaussian mechanism noise. All trajectory parameters are fixed a
priori so the ablation isolates the value of trajectory smoothing.

\subsection{Fixed-mask DCT-Gaussian backend}

Given a mask sequence \(M=(M_1,\ldots,M_T)\), each selected block is converted
to YCbCr. The backend keeps only the \(2\times2\) low-frequency DCT
coefficients of the Y channel. For selected block \(b\) in frame \(t\), the
query vector is clipped as
\begin{equation}
\bar q_{t,b}(Z)=q_{t,b}(Z)\min\{1,C/\|q_{t,b}(Z)\|_2\},
\end{equation}
with \(C=2\), and released as
\begin{equation}
\tilde q_{t,b}=\bar q_{t,b}(Z)+\mathcal{N}(0,\sigma^2I).
\end{equation}
The selected block is reconstructed from the noisy low-frequency Y-DCT
coefficients, zero high-frequency Y coefficients, and neutral chroma. Unselected
blocks are copied from the backend input. This simple backend makes the
protected query explicit and keeps the release transformation fixed across mask
policies.

\subsection{Mask policies}

All policy rows use the same backend and differ only in \(M\). RAT fusion
selects the top-\(k\) fused scores. VideoMAE-only sets the first cue weight to
one and the other weights to zero. RAT+traj. applies the trajectory smoothing
above. Random uses content-independent random scores. Shifted fusion
rolls the RAT-fusion score map by a fixed public spatial offset of three block
rows and four block columns, with wrap-around and the same offset in every
frame. This preserves the score distribution while breaking spatial alignment,
so it tests whether gains come from alignment rather than score concentration.
The Motion reference selects top-\(k\) blocks from frame-difference scores on
\(Z\) and is evaluated against the closely related raw-video proxy \(R(X)\). It
is a deliberately proxy-aligned sanity-check reference.

\section{Fixed-Mask Selected-Query Accounting}
\label{sec:privacy}

\subsection{Threat model and scope}

The protected object is the concatenated vector of clipped \(2\times2\)
low-frequency Y-DCT coefficients from selected blocks. Let \(M\) be fixed
independently of the private backend input, or supplied externally. Define
\(Z\sim_M Z'\) when two backend-input videos agree outside the selected
coordinates of \(M\), and may differ only in selected blocks. The accounting
guarantee is conditioned on \(M\): it accounts for the selected queries under
the mask-restricted adjacency, while video-dependent RAT selection is outside
this selected-query accounting. Because unselected blocks are copied from
\(Z\), arbitrary changes to private unselected content also prevent a
full-rendered-video DP claim. A private end-to-end system would therefore
require accounting for selection and protection of unselected content.

\subsection{R\'enyi-DP accounting}

For one frame, concatenate the clipped selected-block DCT queries. Under
replacement adjacency, clipping each block vector to \(C\) gives sensitivity
\begin{equation}
\Delta \leq 2C\sqrt{k}.
\end{equation}
The Gaussian mechanism with variance \(\sigma^2\) satisfies, at order
\(\alpha>1\),
\begin{equation}
\rho_\alpha=\frac{\alpha\Delta^2}{2\sigma^2}.
\end{equation}
Composing over \(T\) released frames and converting to
\((\epsilon,\delta)\)-DP gives
\begin{equation}
\epsilon(\alpha)=T\rho_\alpha+\frac{\log(1/\delta)}{\alpha-1}.
\end{equation}
Experiments use \(k=20\), \(C=2\), and \(\delta=10^{-4}\). Across the 20-video
subset, \(T\) ranges from 53 to 379, \(\sigma\) ranges from 82.95 to 222.11,
and the optimized accounting uses \(\alpha=4\), giving
\(\epsilon\in[7.95,8.00]\).

\begin{figure*}[t]
\centering
\small
\textbf{(a)}\quad
\includegraphics[width=0.78\linewidth]{figures/E2_k_sensitivity_deltas.png}\\[-0.5ex]
\textbf{(b)}\quad
\includegraphics[width=0.78\linewidth]{figures/mask_policy_example.png}
\caption{Quantitative budget sensitivity and qualitative mask evidence.
(a) RAT-fusion-minus-control differences for \(k\in\{10,20,30,40\}\);
points are means over 20 paired videos and error bars are 95\% bootstrap
confidence intervals. (b) Raw frame and Motion reference, RAT,
Random, and Shifted-fusion masks for one example.}
\label{fig:vis}
\end{figure*}

\section{Experiments}
\label{sec:exp}

\subsection{Protocol and metrics}

% PDF link target: https://github.com/aronlqc/ICASSP-experiment.git
The code and experiment materials for this study are publicly available at
\mbox{GitHub: ICASSP-experiment}.

We evaluate 20 videos from UCF101~\cite{soomro2012ucf101}: four clips each
from BalanceBeam, Basketball, Biking, Diving, and WalkingWithDog. We use a
fixed balanced 20-video subset, with four clips from each of five classes, and
retain exact clip identities in the run manifest. Within each clip, all rows
reuse the same saved pre-release backend input \(Z\), selected area,
DCT-Gaussian backend, Gaussian calibration, \((\epsilon,\delta)\) accounting
target, and Gaussian mechanism RNG stream. Only the mask policy changes.
Gaussian mechanism noise is distinct from selector randomness, and the primary
comparison uses one fixed mask-construction realization for Random and
RAT+traj.; Sec.~\ref{sec:e1} repeats the stochastic construction across ten
master seeds.

The reference \(R_t\) is an operational motion-salience reference built from
frame differences in the raw video. We
compute grayscale absolute differences, use one-sided differences at the first
and last frames and the average of previous and next differences in interior
frames, average to the \(14\times14\) block grid, robustly normalize using the
2--98\% quantile range, and select the top-\(k\) blocks. Cover, RawErr, and
IncErr are the alignment scores defined in Sec.~\ref{sec:method}. Complement
PSNR is computed outside \(R_t\) using raw-video error.

The primary comparison fixes the selected area at \(k=20\) of 196 blocks,
about 10.16\% of pixels after mapping blocks back to the native video grid.
The Gaussian calibration and RNG stream are matched across policies; selected-
block ordering remains policy-specific. Section~\ref{sec:e2} repeats the
comparison at four selected-area budgets.

\begin{table*}[t]
\caption{Aggregate same-backend mask-isolation results on 20 UCF101 videos.
Only the mask policy changes.}
\label{tab:mask}
\centering
\small
\setlength{\tabcolsep}{3pt}
\begin{tabular}{lccccc}
\toprule
\textbf{Policy} & \textbf{PSNR} & \textbf{Cover} & \textbf{RawErr} & \textbf{IncErr} & \textbf{Comp. PSNR} \\
\midrule
Motion reference & 14.8115 & 0.7316 & 0.6727 & 0.7223 & 19.2696 \\
RAT fusion & 14.7087 & 0.2550 & 0.2408 & 0.2452 & 15.3111 \\
RAT+traj. & 14.6923 & 0.2511 & 0.2374 & 0.2415 & 15.2809 \\
VideoMAE-only & 14.3719 & 0.0697 & 0.0730 & 0.0654 & 14.0997 \\
Random & 14.6338 & 0.1030 & 0.1048 & 0.0966 & 14.4951 \\
Shifted fusion & 14.7107 & 0.0832 & 0.0887 & 0.0803 & 14.4885 \\
\bottomrule
\end{tabular}
\par\medskip
\caption{Paired RAT-fusion-minus-control differences over 20 videos. Exact sign
tests give \(p=2.0{\times}10^{-6}\) for Random/Shifted,
\(p\leq2.58{\times}10^{-3}\) for VideoMAE-only, and \(p\geq0.824\) for
RAT+traj.}
\label{tab:paired}
\centering
\small
\setlength{\tabcolsep}{4pt}
\begin{tabular}{lcccc}
\toprule
\textbf{Control} & \textbf{\(\Delta\)Cover} & \textbf{\(\Delta\)RawErr} & \textbf{\(\Delta\)IncErr} & \textbf{Outcome} \\
\midrule
Random & 0.1520 & 0.1360 & 0.1486 & 20/20 on all \\
Shifted fusion & 0.1718 & 0.1522 & 0.1649 & 20/20 on all \\
VideoMAE-only & 0.1853 & 0.1678 & 0.1798 & 17/20--18/20 \\
RAT+traj. & 0.0039 & 0.0034 & 0.0037 & n.s. \\
\bottomrule
\end{tabular}
\end{table*}

\subsection{Mask-isolation results}

Mask policy alone produces a large and systematic shift in where perturbation
lands. At \(k=20\), RAT beats Random on Cover, RawErr, and IncErr for all 20/20
videos, with RAT-minus-Random differences of \(0.1520\), \(0.1360\), and
\(0.1486\). RAT also beats Shifted fusion on all three metrics for all 20/20
videos, with differences of \(0.1718\), \(0.1522\), and \(0.1649\). Exact sign
tests give \(p=2.0{\times}10^{-6}\) for both comparisons. Because Shifted fusion
preserves the score distribution while breaking spatial alignment, its failure
attributes the RAT gain to spatial correspondence rather than score
concentration.

To make the unit of analysis explicit, we re-used the saved per-video logs and
recomputed paired RAT-minus-control differences without regenerating videos.
Figure~\ref{fig:per_video_reanalysis} reports the mean difference with a 95\%
paired bootstrap confidence interval, the median difference, and paired Hedges
$g_z$ for each metric. Against Random, the mean differences are $0.1520$,
$0.1360$, and $0.1486$ for Cover, RawErr, and IncErr, with corresponding
medians $0.1371$, $0.1338$, and $0.1407$ and $g_z=1.47$, $1.42$, and $1.48$.
Against Shifted fusion, the mean differences are $0.1718$,
$0.1522$, and $0.1649$, with medians $0.1714$, $0.1491$, and $0.1668$ and
$g_z=1.58$, $1.50$, and $1.55$. The 95\% bootstrap intervals exclude zero for
both controls on all three metrics. The VideoMAE-only comparison also remains
large (mean IncErr difference $0.1798$, median $0.1687$, $g_z=1.12$), whereas
RAT versus RAT+traj. is small (mean IncErr difference $0.0037$, median
$0.0072$, $g_z=0.10$; 95\% CI $[-0.0119,0.0197]$), confirming that the
trajectory-smoothed variant adds no measurable allocation gain.

\begin{figure*}[t]
\centering
\includegraphics[width=0.98\textwidth]{figures/per_video_effects.pdf}
\caption{Per-video paired reanalysis on the 20-video subset. Circles show mean
RAT-minus-control differences with 95\% percentile paired-bootstrap confidence
intervals; diamonds show median differences; labels report paired Hedges
$g_z$. The analysis reuses the saved per-video metrics and does not regenerate
videos.}
\label{fig:per_video_reanalysis}
\vspace{-0.08in}
\end{figure*}

RAT substantially exceeds VideoMAE-only on 17/20--18/20 videos
\((p\leq2.58{\times}10^{-3})\), showing that semantic salience alone does not
explain the gain. RAT and RAT+traj. are statistically indistinguishable
\((p\geq0.824)\), an informative negative result: additional trajectory
smoothing is unnecessary for the observed allocation gain. The proxy-aligned
Motion reference is higher on all 20 videos because it is deliberately aligned
with the evaluation proxy; it serves as a sanity-check reference. Figure~\ref{fig:vis}(b)
illustrates these masks.

\subsection{E1: Robustness across stochastic realizations}
\label{sec:e1}

The allocation advantage replicates across all ten stochastic realizations.
Across ten master seeds (1,400 successful runs), we first average metrics over
seeds within each video and then pair the 20 videos. RAT beats Random and
Shifted on all videos and metrics; the IncErr differences are \(+0.1505\) and
\(+0.1660\) (\(p=2.0{\times}10^{-6}\)). Every run-level RAT-minus-Random
difference is positive, while RAT and RAT+traj. remain indistinguishable
(\(p=0.503\)).

\subsection{E2: Selected-area sensitivity}
\label{sec:e2}

For \(k=10,20,30,40\), we recalibrate \(\sigma\) per clip and budget to keep
\(\epsilon=8,\delta=10^{-4}\) (560 successful runs). Figure~\ref{fig:vis}(a)
shows positive RAT-minus-Random and RAT-minus-Shifted differences throughout.
All comparisons are 20/20 except \(k=10\) versus Shifted, which is 19/20.
The allocation advantage persists across a four-fold range of selected-area
budgets under matched \(\epsilon=8\) recalibration.

\section{Discussion and Conclusion}
\label{sec:discussion}

Perturbation placement materially changes where release error accumulates even
when the backend is identical. RAT exploits this degree of freedom through
multi-cue fusion of semantic, temporal, texture, and flow signals. At \(k=20\),
it beats Random and Shifted on all three metrics for 20/20 videos, with IncErr
improvements of \(+0.1486\) and \(+0.1649\). The advantage replicates across
ten master seeds and persists for \(k\in\{10,20,30,40\}\) under matched
\(\epsilon=8\) recalibration. VideoMAE-only is weaker, while trajectory
smoothing gives no measurable allocation gain. These results establish mask
policy as a controllable design lever, and same-backend isolation provides a
clean way to study this design variable. The present R\'enyi-DP accounting
applies to selected DCT queries under fixed or external masks; extending this
controlled allocation framework to end-to-end private selection and release is
a natural next step.

\clearpage

\begin{thebibliography}{11}

\bibitem{dwork2006dp}
C. Dwork, F. McSherry, K. Nissim, and A. Smith, ``Calibrating noise to
sensitivity in private data analysis,'' in \emph{Proc. TCC}, 2006.

\bibitem{mironov2017rdp}
I. Mironov, ``R\'enyi differential privacy,'' in \emph{Proc. IEEE CSF}, 2017.

\bibitem{schneider2024maskdp}
D. Schneider, S. Sajadmanesh, V. Sehwag, M. S. Sarfraz, R. Stiefelhagen,
L. Lyu, and V. Sharma, ``Activity recognition on avatar-anonymized datasets
with masked differential privacy,'' arXiv:2410.17098, 2024.

\bibitem{nken2025videodprp}
A. T. A. Nken, S. McKeever, P. Corcoran, and I. Ullah, ``Video-DPRP: A
differentially private approach for visual privacy-preserving video human
activity recognition,'' in \emph{Machine Learning and Knowledge Discovery in
Databases, ECML PKDD 2025, Proceedings}, LNCS 16017, pp. 345--362, 2026,
doi:10.1007/978-3-032-06096-9\_20.

\bibitem{fan2023mgm}
D. Fan, J. Wang, S. Liao, Y. Zhu, V. Bhat, H. Santos-Villalobos, R. MV, and
X. Li, ``Motion-guided masking for spatiotemporal representation learning,''
in \emph{Proc. ICCV}, 2023.

\bibitem{li2023stprivacy}
M. Li, X. Xu, H. Fan, P. Zhou, J. Liu, J.-W. Liu, J. Li, J. Keppo,
M. Z. Shou, and S. Yan, ``STPrivacy: Spatio-temporal privacy-preserving action
recognition,'' in \emph{Proc. ICCV}, pp. 5106--5115, 2023.

\bibitem{ilic2024selective}
F. Ilic, H. Zhao, T. Pock, and R. P. Wildes, ``Selective, interpretable, and
motion consistent privacy attribute obfuscation for action recognition,'' in
\emph{Proc. CVPR}, 2024.

\bibitem{aslam2026tokenprune}
N. Aslam, A. Ray, J. B. Haurum, L. Esterle, and K. Nasrollahi, ``From pixels
to privacy: Temporally consistent video anonymization via token pruning for
privacy preserving action recognition,'' in \emph{Proc. CVPR Workshops}, 2026.

\bibitem{tong2022videomae}
Z. Tong, Y. Song, J. Wang, and L. Wang, ``VideoMAE: Masked autoencoders are
data-efficient learners for self-supervised video pre-training,'' in
\emph{Proc. NeurIPS}, 2022.

\bibitem{teed2020raft}
Z. Teed and J. Deng, ``RAFT: Recurrent all-pairs field transforms for optical
flow,'' in \emph{Proc. ECCV}, 2020.

\bibitem{soomro2012ucf101}
K. Soomro, A. R. Zamir, and M. Shah, ``UCF101: A dataset of 101 human actions
classes from videos in the wild,'' CRCV-TR-12-01, 2012.

\end{thebibliography}
